\scanpage{15}\mspage{14}
Si on ne se borne aux familles finies, \unsure{Con.|on a|ben.} encore connexité. En effet, si $X\in M$ est $\ne\{L\}$, soit $A\in X$ \unsure{tels|tel|telle} que $A\ne L$, Alors $X$ et $X_A$ dans $=$ comp \inserted{connexe et}
\[
\lvert X_A\rvert=A
\]
donc si $\deleted{\{x\}}x\in L\setminus A$, on trouve une multistrate en \unsure{prenant|posant}
\[
Z=\bigl\{B\unsure{,|;}B\cup\{x\},\{x\}\mid B\in X_A\bigr\}
\]
\margin{\unsure{.}}

\uline{NB} Construction qu'on a refait ici plusieurs fois : Soient $\deleted{X}F\in\deleted{M}\unsure{\mathcal F|\mathscr F}$, $\deleted{Y}G\in\deleted{\unsure{\mathcal G|M}}\unsure{\mathcal F|\mathcal G}$, $\lvert\deleted{\unsure{X}}F\rvert\cap\lvert G\rvert=\emptyset$, posons $F\boxtimes G=\{A\cup B\mid A\in\deleted{\unsure{X}}F,\ B\in G,\ \inserted{\text{$A$ et $B$ pas tous deux $\emptyset$}}\}$, alors $(A\cup B)^{\circ}=A^{\circ}\cup B^{\circ}$, $\partial(A\cup B)=(\partial A\cup B)\cup(A\cup\partial B)$ -- si $F$ et $G$ \deleted{sont} \deleted{des} multistrates i.e. \deleted{\unsure{admet|admettent|alors}} \deleted{\unsure{pour|par|pas}} plus grand éléments \deleted{\unsure{et que si}} $A_0,B_0$ alors $A_0\cup B_0$ est pl grd élément de $F\boxtimes G$, qui est une multistrate.

\deleted{Celle-ci contient \unsure{en|un} \unsure{l}} $F\boxtimes G$ contient $F\unsure{\amalg|\sqcup}G$ comme sous-figure. Donc, si $F$ et $G$ sont des multistrates ``disjointes'' elles sont dans la $=$ comp. connexe car \unsure{disjointes|majorées|dominées} par $F\boxtimes G$. Donc en tous cas, \margin{)} $M$ connexe ($=$ sans \deleted{\unsure{la|les}} limiter aux \deleted{familles} figures ensemblistes finies).
